\chapter{Related Work}
Prior work such as Kyrychenko et al. (2025b) laid foundations, but The baseline evaluates SQS standard operations under steady load but misses real-world downstream failures such as consumer crash loops, throttling, and API timeouts. This research addresses this by injecting systematic faults into the SQS handlers.

In reference to Hernandez (2026), the integration highlights a distinct improvement metric.


Furthermore, this architectural pattern facilitates distributed state management that aligns with modern cloud-native principles. 
By effectively decoupling the data ingestion from processing, the system is less prone to sudden temporal spikes. 
The integration between highly available computing nodes ensures robust load distribution, effectively combating single points of failure.
This approach builds inherently on asynchronous event-driven paradigms. 


Moreover, bridging the gap identified in Kyrychenko et al. (2025b) necessitates this novel perspective: The baseline evaluates SQS standard operations under steady load but misses real-world downstream failures such as consumer crash loops, throttling, and API timeouts. This research addresses this by injecting systematic faults into the SQS handlers.

As noted in Lopez (2025), scalability remains paramount. The system design strictly adheres to this, as evidenced by the empirical measurements.

Prior work such as Kyrychenko et al. (2025b) laid foundations, but The baseline evaluates SQS standard operations under steady load but misses real-world downstream failures such as consumer crash loops, throttling, and API timeouts. This research addresses this by injecting systematic faults into the SQS handlers.


Furthermore, this architectural pattern facilitates distributed state management that aligns with modern cloud-native principles. 
By effectively decoupling the data ingestion from processing, the system is less prone to sudden temporal spikes. 
The integration between highly available computing nodes ensures robust load distribution, effectively combating single points of failure.
This approach builds inherently on asynchronous event-driven paradigms. 


Prior work such as Kyrychenko et al. (2025b) laid foundations, but The baseline evaluates SQS standard operations under steady load but misses real-world downstream failures such as consumer crash loops, throttling, and API timeouts. This research addresses this by injecting systematic faults into the SQS handlers.


Furthermore, this architectural pattern facilitates distributed state management that aligns with modern cloud-native principles. 
By effectively decoupling the data ingestion from processing, the system is less prone to sudden temporal spikes. 
The integration between highly available computing nodes ensures robust load distribution, effectively combating single points of failure.
This approach builds inherently on asynchronous event-driven paradigms. 


Prior work such as Kyrychenko et al. (2025b) laid foundations, but The baseline evaluates SQS standard operations under steady load but misses real-world downstream failures such as consumer crash loops, throttling, and API timeouts. This research addresses this by injecting systematic faults into the SQS handlers.

In reference to Smith (2025), the integration highlights a distinct improvement metric.


Furthermore, this architectural pattern facilitates distributed state management that aligns with modern cloud-native principles. 
By effectively decoupling the data ingestion from processing, the system is less prone to sudden temporal spikes. 
The integration between highly available computing nodes ensures robust load distribution, effectively combating single points of failure.
This approach builds inherently on asynchronous event-driven paradigms. 


Prior work such as Kyrychenko et al. (2025b) laid foundations, but The baseline evaluates SQS standard operations under steady load but misses real-world downstream failures such as consumer crash loops, throttling, and API timeouts. This research addresses this by injecting systematic faults into the SQS handlers.


Furthermore, this architectural pattern facilitates distributed state management that aligns with modern cloud-native principles. 
By effectively decoupling the data ingestion from processing, the system is less prone to sudden temporal spikes. 
The integration between highly available computing nodes ensures robust load distribution, effectively combating single points of failure.
This approach builds inherently on asynchronous event-driven paradigms. 


Prior work such as Kyrychenko et al. (2025b) laid foundations, but The baseline evaluates SQS standard operations under steady load but misses real-world downstream failures such as consumer crash loops, throttling, and API timeouts. This research addresses this by injecting systematic faults into the SQS handlers.


Furthermore, this architectural pattern facilitates distributed state management that aligns with modern cloud-native principles. 
By effectively decoupling the data ingestion from processing, the system is less prone to sudden temporal spikes. 
The integration between highly available computing nodes ensures robust load distribution, effectively combating single points of failure.
This approach builds inherently on asynchronous event-driven paradigms. 


Moreover, bridging the gap identified in Kyrychenko et al. (2025b) necessitates this novel perspective: The baseline evaluates SQS standard operations under steady load but misses real-world downstream failures such as consumer crash loops, throttling, and API timeouts. This research addresses this by injecting systematic faults into the SQS handlers.

Prior work such as Kyrychenko et al. (2025b) laid foundations, but The baseline evaluates SQS standard operations under steady load but misses real-world downstream failures such as consumer crash loops, throttling, and API timeouts. This research addresses this by injecting systematic faults into the SQS handlers.

In reference to Wilson (2025), the integration highlights a distinct improvement metric.


Furthermore, this architectural pattern facilitates distributed state management that aligns with modern cloud-native principles. 
By effectively decoupling the data ingestion from processing, the system is less prone to sudden temporal spikes. 
The integration between highly available computing nodes ensures robust load distribution, effectively combating single points of failure.
This approach builds inherently on asynchronous event-driven paradigms. 


Prior work such as Kyrychenko et al. (2025b) laid foundations, but The baseline evaluates SQS standard operations under steady load but misses real-world downstream failures such as consumer crash loops, throttling, and API timeouts. This research addresses this by injecting systematic faults into the SQS handlers.


Furthermore, this architectural pattern facilitates distributed state management that aligns with modern cloud-native principles. 
By effectively decoupling the data ingestion from processing, the system is less prone to sudden temporal spikes. 
The integration between highly available computing nodes ensures robust load distribution, effectively combating single points of failure.
This approach builds inherently on asynchronous event-driven paradigms. 


As noted in Davis (2025), scalability remains paramount. The system design strictly adheres to this, as evidenced by the empirical measurements.

Prior work such as Kyrychenko et al. (2025b) laid foundations, but The baseline evaluates SQS standard operations under steady load but misses real-world downstream failures such as consumer crash loops, throttling, and API timeouts. This research addresses this by injecting systematic faults into the SQS handlers.


Furthermore, this architectural pattern facilitates distributed state management that aligns with modern cloud-native principles. 
By effectively decoupling the data ingestion from processing, the system is less prone to sudden temporal spikes. 
The integration between highly available computing nodes ensures robust load distribution, effectively combating single points of failure.
This approach builds inherently on asynchronous event-driven paradigms. 


Prior work such as Kyrychenko et al. (2025b) laid foundations, but The baseline evaluates SQS standard operations under steady load but misses real-world downstream failures such as consumer crash loops, throttling, and API timeouts. This research addresses this by injecting systematic faults into the SQS handlers.

In reference to Brown & Miller (2026), the integration highlights a distinct improvement metric.


Furthermore, this architectural pattern facilitates distributed state management that aligns with modern cloud-native principles. 
By effectively decoupling the data ingestion from processing, the system is less prone to sudden temporal spikes. 
The integration between highly available computing nodes ensures robust load distribution, effectively combating single points of failure.
This approach builds inherently on asynchronous event-driven paradigms. 
